% =====================================================================
\chapter*{Instrucciones}
\addcontentsline{toc}{chapter}{Instrucciones}
\markboth{Instrucciones}{}
% =====================================================================

Este examen cierra la Fase~V y el curso. Evalúa cómo un modelo de lenguaje preentrenado se convierte en un asistente: el ajuste supervisado y eficiente en parámetros, el aprendizaje por refuerzo y las preferencias, la recuperación de información y los agentes, y las herramientas con las que se estudia por dentro y se evalúa su seguridad. Las técnicas de las fases anteriores (entropía cruzada, SVD, softmax, intervalos de confianza) se dan por sabidas y se usan sin repetir su desarrollo.

\section*{Condiciones del examen}
\begin{center}
\begin{tabularx}{\textwidth}{@{}lX@{}}
\toprule
Duración & 3 horas. Reparto orientativo: 45 minutos para la parte~A y 135 minutos para la parte~B. \\
Material permitido & Calculadora científica y una hoja A4 de fórmulas escrita a mano por el estudiante. No se permiten apuntes, libros ni dispositivos con conexión. \\
Puntuación total & 100 puntos: parte~A, 30 puntos; parte~B, 70 puntos. \\
Aprobado & 70 puntos o más (70\,\%). \\
Módulos evaluados & M20 a M23 (véase la tabla siguiente). \\
\bottomrule
\end{tabularx}
\end{center}

\section*{Normas}
\begin{enumerate}[label=\arabic*.]
  \item \textbf{Parte A. Cuestiones} (30 puntos): 20 preguntas de opción múltiple, de 1,5 puntos cada una. Cada pregunta tiene una sola respuesta correcta. Las respuestas en blanco y las erróneas valen 0 puntos; las erróneas no restan.
  \item \textbf{Parte B. Problemas} (70 puntos): 7 ejercicios de desarrollo, de 10 puntos cada uno: dos por cada uno de los módulos M20, M21 y M22 y uno del M23. La puntuación de cada apartado figura entre corchetes.
  \item En la parte~B debe escribirse el razonamiento completo. Un resultado correcto sin justificación se puntúa con la mitad de los puntos del apartado.
  \item Los resultados se dan con 3 cifras significativas salvo que se indique otra cosa; en los pasos intermedios conviene conservar 4 o más. Se escribe $\ln$ (o $\log$) para el logaritmo natural, de modo que las pérdidas y las divergencias se miden en nats.
  \item Notación: $\pi_\theta$ es la política (el modelo que se entrena), $\piref$ el modelo de referencia, $\sigma$ la función sigmoide, $\beta$ el coeficiente de la penalización KL, $r$ el rango de LoRA y $s=\alpha/r$ su escala; $1\,\mathrm{GB}=10^{9}$ bytes.
  \item La \emph{Parte C. Proyecto de fase} se entrega aparte, en el plazo de tres semanas, y cuenta en la nota de la fase, no en la de este examen.
\end{enumerate}

\section*{Módulos evaluados}
\begin{center}
\begin{tabular}{@{}cllc@{}}
\toprule
Módulo & Título & Cuestiones & Ejercicios \\
\midrule
M20 & Ajuste fino & 1--5 & B.1 y B.2 \\
M21 & Aprendizaje por refuerzo y preferencias & 6--11 & B.3 y B.4 \\
M22 & Agentes, recuperación y RAG & 12--16 & B.5 y B.6 \\
M23 & Interpretabilidad y seguridad & 17--20 & B.7 \\
\bottomrule
\end{tabular}
\end{center}

\section*{Procedencia de los ejercicios}
Los ejercicios B.1, B.3, B.4, B.5 y B.6 están adaptados, total o parcialmente, de tareas y exámenes de cursos universitarios: la tarea~5 de Stanford CS336 y su suplemento (primavera de 2026), el examen parcial de Stanford CS234 (invierno de 2018), los ejercicios del capítulo~8 de \emph{Introduction to Information Retrieval}, texto del curso CS276 de Stanford, y los problemas de clase de MIT 18.05 (primavera de 2022). La fuente se indica al pie de cada enunciado y en la bibliografía. Los enunciados son de redacción propia y los datos se han cambiado para que puedan resolverse a mano. Las soluciones completas están en el \cref{ap:soluciones}, y la tabla de autoevaluación, en el \cref{ap:autoevaluacion}.
